
\chapter{Сквозной пример: Calculator}
\label{chap:calculator}

Рассмотрим Calculator --- приложение, которое
пошагово вычисляет арифметические выражения
и отображает промежуточные результаты.

Программа состоит из трёх автоматных объектов:

\begin{itemize}
  \item \code{CalculatorController} управляет вычислением;
  \item \code{OperatorService} предоставляет операции
    для определения приоритета и вычисления операторов;
  \item \code{CalculatorView} обрабатывает действия
    пользователя и обновляет интерфейс.
\end{itemize}

Полная реализация находится в
\RepoPath{examples/calculator}.

\section{Как работает Calculator}

Вычисление состоит из двух основных этапов:
преобразования выражения в обратную польскую
нотацию (ОПН) и вычисления полученной записи.

Например, выражение \code{2 + 3 * 4} сначала
преобразуется в \code{2 3 4 * +}.

В ОПН порядок операций определяется расположением
элементов, поэтому скобки и правила приоритета
больше не нужны.

Затем выражение вычисляется с помощью стека:

\begin{table}[ht]
  \centering
  \caption{Вычисление выражения в ОПН}
  \begin{tabularx}{0.72\textwidth}{@{}cX@{}}
    \toprule
    Токен & Содержимое стека после обработки \\
    \midrule
    \code{2} & \code{[2]} \\
    \code{3} & \code{[2, 3]} \\
    \code{4} & \code{[2, 3, 4]} \\
    \code{*} & \code{[2, 12]} \\
    \code{+} & \code{[14]} \\
    \bottomrule
  \end{tabularx}
\end{table}

Числа помещаются в стек, а операторы извлекают
необходимые аргументы и возвращают результат
обратно в стек.

Controller выполняет эти этапы постепенно.
Каждая команда \code{tick} обрабатывает один токен
либо завершает текущую фазу вычисления.

\section{Взаимодействие автоматов}

Во время преобразования выражения Controller
использует Query для получения приоритета операторов
от \code{OperatorService}.

При вычислении ОПН он обращается к тому же сервису
за результатом арифметической операции.

Оба запроса выполняются синхронно, поэтому Controller
может сразу использовать полученные значения.

Для отображения промежуточных результатов
Controller отправляет View объект
\code{CalculatorSnapshot} с данными вычисления.

View обрабатывает полученные данные, обновляет
представление и отправляет команду
\code{draw_completed} в качестве подтверждения.

После подтверждения Controller может продолжить
вычисление или перейти к следующей фазе.

Так вычисление и отображение согласуются через
команды, не требуя прямого обращения Controller
к элементам интерфейса.

\section{Сборка программы}

Рассмотрим основные соединения Calculator:

\lstinputlisting[
  basicstyle=\ttfamily\footnotesize,
  caption={Основные соединения Calculator},
  label={lst:calculator}
]{code/calculator_program.py}

Здесь используются оба механизма взаимодействия:
Query для получения результатов от
\code{OperatorService} и Command для обмена
сообщениями между Controller и View.

Полная сборка программы также включает команды
управления вычислением: \code{tick},
\code{reset}, \code{finish} и подтверждение
\code{draw_completed}.

\section{Режимы работы}

Calculator поддерживает пошаговый и автоматический
режимы вычисления.

В пошаговом режиме пользователь нажимает кнопку
«Следующий шаг», чтобы обработать очередной токен.
После обновления представления кнопка снова
становится доступной.

В автоматическом режиме Controller самостоятельно
планирует следующие команды. При этом вычисление
по-прежнему выполняется отдельными шагами Runtime,
а интерфейс обновляется между ними.

Режим запуска передаётся через \code{StartInput}.
Controller не обращается к Window и не работает
напрямую с Tkinter-виджетами.

\begin{trybox}
Запустите Calculator с выражением \code{2 + 3 * 4}
в пошаговом режиме.

Проследите, как формируется ОПН и изменяется
содержимое стека. Затем попробуйте выражение
\code{(2 + 3) * 4} и сравните результаты.

После этого запустите вычисление в автоматическом
режиме.
\end{trybox}
